% Fig 3 - Predictor detail.
%
% Zoom-in on the 6-layer AR transformer predictor used in JEPA.
% Two explicit blocks shown, an ellipsis in between, and the last block.
% Each block exposes:
%   - Multi-head self-attention (16 heads) with RoPE on Q,K and a causal mask
%   - FFN with AdaLN-Zero modulation conditioned on c = (G, D, Y)
% Inputs: sequence z_{t-T+1} ... z_t, time embeddings dt, conditioning c.
% Output: zhat_{t+1}.
%
% Standalone-compilable: pdflatex fig3_predictor_detail.tex

\documentclass[border=10pt,tikz]{standalone}
\usepackage{tikz}
\usepackage{amsmath,amssymb}
\usetikzlibrary{positioning,arrows.meta,fit,shapes.geometric,calc,backgrounds,patterns}

\begin{document}

\begin{tikzpicture}[
    font=\small,
    box/.style={draw, rounded corners=2pt, align=center, inner sep=5pt, thick},
    submod/.style={draw, rounded corners=2pt, fill=orange!18, align=center,
                   inner sep=4pt, font=\footnotesize, minimum width=58mm,
                   minimum height=10mm},
    state/.style={draw, rectangle, rounded corners=1pt, fill=white,
                  minimum height=10mm, minimum width=44mm, inner sep=3pt},
    cond/.style={box, fill=yellow!18, minimum width=36mm, minimum height=12mm,
                 font=\footnotesize},
    note/.style={align=center, font=\footnotesize\itshape, text=black!60},
    arr/.style={-{Stealth[length=2.5mm]}, thick},
    arrshort/.style={-{Stealth[length=2mm]}, semithick},
    arrcond/.style={-{Stealth[length=2mm]}, semithick, yellow!50!black, dashed}
  ]

  % ========= input sequence =========
  \node[state] (zseq) {$\{\mathbf{z}_{t-T+1},\, \ldots,\, \mathbf{z}_{t}\}$ \\[1pt]
                       \footnotesize sequence length $T$};

  % time embeddings to the left
  \node[state, left=14mm of zseq, fill=blue!4] (dtemb)
    {$\Delta t$ embeddings \\[1pt] \footnotesize RoPE positions};

  \draw[arr] (dtemb) -- (zseq);

  % conditioning c to the right
  \node[cond, right=24mm of zseq] (cond)
    {$\mathbf{c} = (G, D, Y)$ \\[1pt] gust parameters};

  % ========= Block 1 (explicit, two submodules with ample spacing) =========
  % Place attention then FFN, then fit a frame around them.
  \node[submod, below=16mm of zseq] (b1a)
    {Multi-head self-attention (16 heads, $d_{\text{model}}{=}384$) \\
     RoPE on $Q,K$; causal mask};
  \node[submod, below=4mm of b1a] (b1f)
    {FFN (hidden $1536$, dropout $0.1$) \\
     AdaLN-Zero modulation by $\mathbf{c}$};
  \draw[arrshort] (b1a.south) -- (b1f.north);

  % ========= Block 2 (explicit) =========
  \node[submod, below=16mm of b1f] (b2a)
    {Multi-head self-attention (16 heads) \\ RoPE on $Q,K$; causal mask};
  \node[submod, below=4mm of b2a] (b2f)
    {FFN (hidden $1536$, dropout $0.1$) \\ AdaLN-Zero modulation by $\mathbf{c}$};
  \draw[arrshort] (b2a.south) -- (b2f.north);

  % ========= ellipsis =========
  \node[below=8mm of b2f, font=\Large] (dots) {$\vdots$};

  % ========= Block 6 (last, explicit) =========
  \node[submod, below=8mm of dots] (b6a)
    {Multi-head self-attention (16 heads) \\ RoPE on $Q,K$; causal mask};
  \node[submod, below=4mm of b6a] (b6f)
    {FFN (hidden $1536$, dropout $0.1$) \\ AdaLN-Zero modulation by $\mathbf{c}$};
  \draw[arrshort] (b6a.south) -- (b6f.north);

  % ========= output =========
  \node[state, below=14mm of b6f, fill=orange!4] (zout)
    {$\widehat{\mathbf{z}}_{t+1} \in \mathbb{R}^{d}$};

  % ========= main forward arrows (between blocks) =========
  \draw[arr] (zseq.south) -- (b1a.north);
  \draw[arr] (b1f.south)  -- (b2a.north);
  \draw[arr] (b2f.south)  -- (dots.north);
  \draw[arr] (dots.south) -- (b6a.north);
  \draw[arr] (b6f.south)  -- (zout.north);

  % ========= block frames (fit each block around its two submodules) =========
  \begin{scope}[on background layer]
    \node[draw=orange!50, rounded corners=3pt, fill=orange!6, inner sep=4pt,
          fit=(b1a)(b1f),
          label={[orange!60!black,font=\footnotesize\bfseries]left:Block 1}] (b1) {};
    \node[draw=orange!50, rounded corners=3pt, fill=orange!6, inner sep=4pt,
          fit=(b2a)(b2f),
          label={[orange!60!black,font=\footnotesize\bfseries]left:Block 2}] (b2) {};
    \node[draw=orange!50, rounded corners=3pt, fill=orange!6, inner sep=4pt,
          fit=(b6a)(b6f),
          label={[orange!60!black,font=\footnotesize\bfseries]left:Block 6}] (b6) {};
  \end{scope}

  % ========= conditioning arrows (yellow dashed) -- to each block =========
  \draw[arrcond] (cond.south) to[out=-90,in=0] (b1.east);
  \draw[arrcond] (cond.south) to[out=-90,in=0] (b2.east);
  \draw[arrcond] (cond.south) to[out=-90,in=0] (b6.east);

  % ========= causal mask indicator on the right side =========
  % Placed FAR right so the conditioning arrows from c to each block route
  % cleanly along the right edge of the spine without crossing the mask widget
  % or the hyperparameter summary.
  \node[draw, fill=white, inner sep=2pt, right=44mm of b2,
        font=\footnotesize, align=center] (mask)
    {Causal mask \\[2pt]
     \begin{tikzpicture}[baseline=0pt]
       % 4x4 lower-triangular indicator
       \foreach \r in {0,1,2,3} {
         \foreach \c in {0,1,2,3} {
           \ifnum\c>\r
             \fill[gray!20] (\c*3mm, -\r*3mm) rectangle ++(3mm, -3mm);
           \else
             \fill[orange!55!black!70] (\c*3mm, -\r*3mm) rectangle ++(3mm, -3mm);
           \fi
           \draw[gray!50, very thin] (\c*3mm, -\r*3mm) rectangle ++(3mm, -3mm);
         }
       }
     \end{tikzpicture} \\[2pt]
     future $\to$ past blocked};

  % ========= hyperparameter summary box =========
  \node[note, below=18mm of mask, draw=orange!40, dashed,
        rounded corners=2pt, inner sep=4pt, align=left] (hp)
    {\textbf{Hyperparameters} \\
     6 layers, hidden $384$ \\
     16 attention heads \\
     dropout $0.1$ \\
     AdamW, lr $5{\times}10^{-4}$};

  % ========= background frame =========
  \begin{scope}[on background layer]
    \node[draw=orange!50, rounded corners=4pt, dashed, inner sep=6mm,
          fit=(dtemb)(zseq)(cond)(b1)(b2)(dots)(b6)(zout)(mask)(hp),
          label={[orange!60!black,font=\small]above:Autoregressive transformer predictor $P_{\phi}$ (causal, 6 blocks)}] {};
  \end{scope}

\end{tikzpicture}

\end{document}
